\documentclass{article}
\usepackage[UTF8]{ctex}
\usepackage{amsmath}
\usepackage{amssymb}
\usepackage{graphicx}
\usepackage{geometry}
\geometry{a4paper, margin=2cm}

\newcommand{\problemname}[1]{\begin{center}\Large\bfseries #1\end{center}}
\usepackage{enumitem}

\begin{document}
\problemname{求面积的问题}

小蒜的老师准备了很多个口袋，每个口袋里面有两张纸条。

\begin{itemize}
  \item 第一张纸条上说明了我们将要面对的是一个怎样的图形（可能是三角形，也可能是矩形）

  \item 第二张纸条上说明了两条边的长度 $m$ 和 $n$（分别对应三角形的底和高，或者矩形的长和宽）

\end{itemize}
小蒜会从口袋中随机挑一个，并把纸条上的信息给你，让你的程序计算出图形的面积。

提示

\begin{itemize}
  \item 矩形面积：长 $*$ 宽；

  \item 三角形面积：低 $*$ 高 $/\text{ }2$；

\end{itemize}
\subsection*{输入格式}
输入包含两行。

输入的第一行为一个字符，只可能是 \texttt{r} 或 \texttt{t}。如果是 \texttt{r} 则表示我们要计算面积的图形是一个矩形；如果是 \texttt{t} 则为三角形。

输入的第二行为两个实数，$m$ 和 $n$。分别对应三角形的底和高（或对应矩形的长和宽）（$0.0<m\le100.0$）。

\subsection*{输出格式}
输出为一行，为目标图形的面积。

\subsection*{样例输入}
\begin{verbatim}
r
3.0 2.1
\end{verbatim}

\subsection*{样例输出}
\begin{verbatim}
6.3
\end{verbatim}

\subsection*{样例 2 输入}
\begin{verbatim}
t
3.0 5.0
\end{verbatim}

\subsection*{样例输出}
\begin{verbatim}
7.5
\end{verbatim}


\end{document}
